%%%%%%%%%%%%%%%%%%%%%%%%%%%%%%%%%
%
%       EXERCÍCIO
%
%%%%%%%%%%%%%%%%%%%%%%%%%%%%%%%%%

\ifdefstring{\atividade}{prova}{%
    \ifdefstring{\modo}{objetivo}{%
        \renewcommand{\valorquestao}{\ValorQObj\ ponto}
    }{%
        \renewcommand{\valorquestao}{\ValorQDisc\ pontos}
    }%
}%

\begin{exercicioBanco}[\valorquestao]
Considere a função
%
\[
    f(x)= \left\{
    \begin{array}{cl}
        x+2, & \text{se } x\le -1,\\[2pt]
        1, & \text{se } -1<x<1,\\[2pt]
        -2x+3, & \text{se } x\ge 1.
    \end{array} 
    \right.
\]

Assinale a alternativa correta.

% Define as alternativas
\newcommand{\alternativas}{%
    \begin{center}
        \begin{tabularx}{\textwidth}{XX}
            (a) A função é crescente em todo \(\bbR\). &
            (b) A função possui exatamente uma raiz real. \\[5pt]
            (c) A imagem da função é \(\bbR\). & 
            (d) A função possui valor máximo igual a \(1\). \\[5pt]
            (e) A função é injetiva.
        \end{tabularx}
    \end{center}
}

% Define a resposta correta
\newcommand{\resposta}{D}

% Lógica condicional para exibição
\ifdefstring{\atividade}{lista}{%
    \alternativas
    \vspace{0.5em}
    
    \noindent\textbf{Resposta:} letra \textbf{\resposta}.
}{%
    \ifdefstring{\modo}{objetivo}{%
        \alternativas
    }{%
        % modo = discursiva → não mostra alternativas
    }
}
\end{exercicioBanco}

